\subsection{トピックと参考文献の対応}
この表は、SWEBOK が示す「トピック対参考資料の行列」を、学習用に簡略化したものである。詳しく学ぶときは、各トピックに対応する文献や標準を読む。

\paragraph{表11：トピックと主な参考資料}
\par\smallskip\noindent\refstepcounter{table}\label{tab:ch10-11}表 \thetable: 表11：トピックと主な参考資料におけるトピック・主な参考資料\par\smallskip
表 \thetable は、表11：トピックと主な参考資料におけるトピック・主な参考資料を比較・確認しやすい形で整理したものである。\par\smallskip
\begin{longtable}{>{\raggedright\arraybackslash}p{0.28\linewidth}>{\raggedright\arraybackslash}p{0.64\linewidth}}
\toprule
トピック & 主な参考資料\\
\midrule
\endfirsthead
\toprule
トピック & 主な参考資料\\
\midrule
\endhead
1. ソフトウェア工学プロセスの基礎 & ISO/IEC/IEEE 12207、PMBOK Guide、ISO/IEC/IEEE 24765、ISO/IEC 25000、ISO/IEC/IEEE 24748。\\
1.1 導入 & ISO/IEC/IEEE 12207、PMBOK Guide。\\
1.2 プロセス定義 & ISO/IEC/IEEE 12207、ISO/IEC/IEEE 24765、ISO/IEC 25000、ISO/IEC/IEEE 24748、ISO/IEC TR 29110。\\
2. ライフサイクル & Sommerville、Farley、PMI Agile Practice Guide、ISO/IEC/IEEE 32675、ISO/IEC/IEEE 24774。\\
2.1 定義・分類・用語 & ISO/IEC/IEEE 12207、Sommerville、Farley、PMBOK Guide。\\
2.2 ライフサイクルが必要な理由 & Farley、ISO/IEC/IEEE 24774。\\
2.3 プロセスモデルとライフサイクルモデル & Sommerville、PMI Agile Practice Guide、ISO/IEC/IEEE 24765。\\
2.4 開発ライフサイクルモデルの考え方 & Sommerville、Farley、Shore and Warden、PMI Agile Practice Guide、ISO/IEC/IEEE 24765、ISO/IEC/IEEE 32675。\\
2.5 開発ライフサイクルモデルと工学的側面 & Sommerville、Farley、Shore and Warden、PMI Agile Practice Guide、Agile Manifesto、McConnell、Agile Alliance、Eckstein and Buck、RUP、OpenUP。\\
2.6 SLCPの管理 & ISO/IEC/IEEE 24748。\\
2.7 プロセス管理 & ISO/IEC/IEEE 12207、ISO/IEC/IEEE 24765。\\
2.8 ライフサイクル適応 & PMI Software Extension、ISO/IEC/IEEE 24748、ISO/IEC 29110。\\
2.10 プロセス基盤・ツール・手法 & Sommerville、Farley、ISO/IEC/IEEE 24765。\\
2.11 プロセス監視 & ISO/IEC/IEEE 12207、Sommerville、Laporte and April、Farley。\\
3. プロセス評価と改善 & Laporte and April、Shewhart and Deming、ISO/IEC 33001、CMMI、Fenton and Bieman。\\
3.4 アジャイルの評価と改善 & Shore and Warden、Dingsøyr。\\
\bottomrule
\end{longtable}
